\section{Summary}\label{sec:ch4_preliminaries_summary}

In this chapter, we established the formalisms used for the remainder of the dissertation. Building on the conceptual account of novelty in \cref{ch:novelty_open_world}, we defined open-world environments as settings in which an agent's knowledge is incomplete with respect to the true state of the world, introduced the crafting running example, and clarified the relationship between environment states, symbolic state descriptions, and subsymbolic state representations.

We then developed the agent framework, integrating symbolic planning (\cref{sec:symbolic_planning}), reinforcement learning (\cref{sec:rl}), and the integrated planning task (IPT) formalism (\cref{sec:hybrid}) that couples them through a detector $\Abstraction : \States \rightarrow \SymStates$ and an executor mapping $e : \Operators \rightarrow \Executors$. The framework distinguishes planning success (plannability) from execution success (solvability).

Within that framework, we formalized novelty as agent-relative mismatch (\cref{def:novelty}) and introduced complementary taxonomies by task impact---prohibitive, obstructive, beneficial, and nuisance---and by domain element---object, attribute, and action. The IPT makes explicit where novelty can manifest: at the symbolic level through incomplete models (prohibitive), or at the execution level through failed executors while symbolic plans remain valid (obstructive).

Finally, we introduced the concept of structured abstractions and argued that the choice of abstraction determines whether novelty-induced mismatch leads to brittle failure or can be managed through efficient adaptation and generalizable behavior. Action abstractions expose the compositional structure of plans, enabling failure localization and targeted learning. Interaction abstractions capture the physically meaningful invariants of manipulation tasks, enabling transfer across objects and embodiments.

The subsequent chapters instantiate this framework. In \cref{ch:benchmarks}, we present benchmark environments and evaluation methodologies that operationalize the novelty taxonomy through controlled injection of object, attribute, and action novelties and provide metrics for measuring adaptation and recovery. In \cref{ch:rapid_learn}, we present hybrid planning--learning methods that exploit action abstractions to achieve efficient adaptation under novelty. In \cref{ch:flex,ch:pho}, we present force-based learning approaches that use interaction abstractions to enable generalizable robotic manipulation. In each case, chapter-specific related work is discussed in the context of the particular contribution.
